% !TEX program = pdflatex
% Plan detaille du papier Legal Knowledge Graph -- 2026-07-28
% Support de discussion avant redaction du manuscrit.

\documentclass[aspectratio=169,10pt]{beamer}

\usepackage[utf8]{inputenc}
\usepackage[T1]{fontenc}
\usepackage[french]{babel}
\usepackage{lmodern}
\usepackage{booktabs}
\usepackage{array}
\usepackage{graphicx}
\usepackage[table]{xcolor}
\newcolumntype{L}[1]{>{\raggedright\arraybackslash}p{#1}}

\usetheme{Madrid}
\usecolortheme{seahorse}
\setbeamertemplate{navigation symbols}{}
\setbeamertemplate{footline}[frame number]
\setbeamerfont{footnote}{size=\tiny}

\definecolor{accent}{RGB}{41, 92, 155}
\definecolor{accentsoft}{RGB}{230, 238, 247}
\definecolor{ok}{RGB}{30, 130, 70}
\definecolor{crosscol}{RGB}{110, 60, 130}
\definecolor{warning}{RGB}{190, 95, 0}

\newcommand{\keybox}[1]{%
  \begin{center}
  \colorbox{accentsoft}{\parbox{0.92\linewidth}{\centering #1}}
  \end{center}}
\newcommand{\warnbox}[1]{%
  \begin{center}
  \colorbox{red!10}{\parbox{0.92\linewidth}{\centering #1}}
  \end{center}}
\newcommand{\decisionbox}[1]{%
  \begin{center}
  \colorbox{orange!14}{\parbox{0.92\linewidth}{\centering\color{warning} #1}}
  \end{center}}
\newcommand{\validated}{\textcolor{ok}{\textbf{Validé}}}
\newcommand{\conditional}{\textcolor{red!75!black}{\textbf{Conditionnel}}}
\newcommand{\todecide}{\textcolor{warning}{\textbf{À décider}}}
\newcommand{\papersection}[1]{%
  {\scriptsize\textcolor{accent}{\textbf{Section du papier -- #1}}}\par\vspace{0.55em}}

\title[Plan papier KG juridique FR]{Construction d'un Knowledge Graph juridique français}
\subtitle{Un plan guidé par les questions du reviewer}
\author{Matthieu Kaeppelin}
\institute{Stage FE Recherche}
\date{28 juillet 2026}

\begin{document}

% 1
\begin{frame}
  \titlepage
\end{frame}

% 2
\begin{frame}{Comment lire ce plan ?}
  \papersection{Vue d'ensemble}
  \footnotesize
  \begin{columns}[c,onlytextwidth]
    \begin{column}{0.27\textwidth}
      \centering
      \colorbox{accentsoft}{\parbox{0.90\linewidth}{\centering\vspace{1.0em}
      \textbf{Question du reviewer}\\[0.45em]
      Ce qu'il doit comprendre
      \vspace{1.0em}}}
    \end{column}
    \begin{column}{0.06\textwidth}\centering$\Longrightarrow$\end{column}
    \begin{column}{0.27\textwidth}
      \centering
      \colorbox{green!10}{\parbox{0.90\linewidth}{\centering\vspace{1.0em}
      \textbf{Réponse du projet}\\[0.45em]
      Ce que nous pouvons défendre
      \vspace{1.0em}}}
    \end{column}
    \begin{column}{0.06\textwidth}\centering$\Longrightarrow$\end{column}
    \begin{column}{0.27\textwidth}
      \centering
      \colorbox{orange!12}{\parbox{0.90\linewidth}{\centering\vspace{1.0em}
      \textbf{Section du papier}\\[0.45em]
      Où la preuve sera portée
      \vspace{1.0em}}}
    \end{column}
  \end{columns}

  \vspace{1.0em}
  \keybox{\scriptsize Le fil du deck suit les questions proposées par les guides d'écriture CVPR,
  puis les relie au plan concret de notre manuscrit.}

  {\tiny Sources méthodologiques : Lepetit, pp.~29--30 ; Freeman, pp.~15, 18 et 25--29.}
\end{frame}

% 3
\begin{frame}{Quelles sections devront répondre à ces questions ?}
  \papersection{Vue d'ensemble du manuscrit}
  \footnotesize
  \begin{columns}[T,onlytextwidth]
    \begin{column}{0.48\textwidth}
      \begin{enumerate}\setlength{\itemsep}{0.55em}
        \item \textbf{Introduction}\\importance, difficulté, idée et contributions ;
        \item \textbf{Related Work}\\solutions existantes et manque précis ;
        \item \textbf{Task and Benchmark}\\sorties, données, golds et périmètre ;
        \item \textbf{Method}\\graphes, retrieval et LLM reranking.
      \end{enumerate}
    \end{column}
    \begin{column}{0.48\textwidth}
      \begin{enumerate}\setcounter{enumi}{4}\setlength{\itemsep}{0.55em}
        \item \textbf{Experiments}\\protocole, baselines et ablations ;
        \item \textbf{Results}\\preuves, diagnostics et résultats négatifs ;
        \item \textbf{Discussion and Limitations}\\portée et menaces à la validité ;
        \item \textbf{Conclusion}\\message soutenu par les preuves finales.
      \end{enumerate}
    \end{column}
  \end{columns}
\end{frame}

% 4
\begin{frame}{Pourquoi ce problème est-il important ?}
  \papersection{1. Introduction -- Motivation}
  \footnotesize
  \textbf{Un avocat, un professionnel du droit ou un particulier cherche à répondre à une question juridique.}

  \vspace{0.65em}
  \begin{columns}[T,onlytextwidth]
    \begin{column}{0.46\textwidth}
      \textbf{\color{accent}1. Les articles : la règle}
      \begin{itemize}\setlength{\itemsep}{0.38em}
        \item Ils définissent le cadre normatif applicable.
        \item Ils indiquent les conditions, exceptions et sanctions prévues par le droit.
      \end{itemize}
    \end{column}
    \begin{column}{0.46\textwidth}
      \textbf{\color{ok}2. Les jurisprudences : l'application}
      \begin{itemize}\setlength{\itemsep}{0.38em}
        \item Elles montrent comment les règles sont interprétées.
        \item Elles relient le texte à des situations concrètes et à des arguments précis.
      \end{itemize}
    \end{column}
  \end{columns}

  \vspace{0.85em}
  \keybox{\scriptsize Répondre correctement suppose de retrouver les deux fondements :
  la règle de droit et son application jurisprudentielle.}
\end{frame}

% 5
\begin{frame}{Pourquoi ce problème est-il difficile ?}
  \papersection{1. Introduction -- Difficultés}
  \scriptsize
  \begin{columns}[T,onlytextwidth]
    \begin{column}{0.23\textwidth}
      \textbf{\color{accent}Expertise juridique}
      \begin{itemize}\setlength{\itemsep}{0.3em}
        \item vocabulaire spécialisé ;
        \item qualifications dépendantes du contexte ;
        \item interprétation experte nécessaire.
      \end{itemize}
    \end{column}
    \begin{column}{0.23\textwidth}
      \textbf{\color{ok}Échelle et cas particuliers}
      \begin{itemize}\setlength{\itemsep}{0.3em}
        \item grand volume de textes et décisions ;
        \item exceptions nombreuses ;
        \item configurations factuelles variées.
      \end{itemize}
    \end{column}
    \begin{column}{0.23\textwidth}
      \textbf{\color{crosscol}Organisation du droit}
      \begin{itemize}\setlength{\itemsep}{0.3em}
        \item domaines et sous-domaines ;
        \item hiérarchies et renvois ;
        \item citations et évolutions temporelles.
      \end{itemize}
    \end{column}
    \begin{column}{0.23\textwidth}
      \textbf{\color{warning}Limites de la similarité}
      \begin{itemize}\setlength{\itemsep}{0.3em}
        \item termes proches, solutions opposées ;
        \item régimes distincts, vocabulaire commun ;
        \item relations absentes du texte local.
      \end{itemize}
    \end{column}
  \end{columns}

  \vspace{0.65em}
  \keybox{\scriptsize La proximité sémantique est utile, mais elle ne suffit pas à représenter
  les oppositions juridiques, les hiérarchies et les liens entre sources.}
\end{frame}

% 6
\begin{frame}{Pourquoi le retrieval est-il indispensable aux assistants juridiques ?}
  \papersection{1. Introduction -- Motivation système}
  \footnotesize
  \begin{columns}[c,onlytextwidth]
    \begin{column}{0.19\textwidth}\centering
      \colorbox{accentsoft}{\parbox{0.90\linewidth}{\centering\vspace{0.65em}\textbf{Question juridique}\vspace{0.65em}}}
    \end{column}
    \begin{column}{0.05\textwidth}\centering$\Longrightarrow$\end{column}
    \begin{column}{0.20\textwidth}\centering
      \colorbox{accentsoft}{\parbox{0.90\linewidth}{\centering\vspace{0.65em}\textbf{Retriever}\\sources candidates\vspace{0.65em}}}
    \end{column}
    \begin{column}{0.05\textwidth}\centering$\Longrightarrow$\end{column}
    \begin{column}{0.20\textwidth}\centering
      \colorbox{green!10}{\parbox{0.90\linewidth}{\centering\vspace{0.65em}\textbf{LLM}\\sélection et synthèse\vspace{0.65em}}}
    \end{column}
    \begin{column}{0.05\textwidth}\centering$\Longrightarrow$\end{column}
    \begin{column}{0.20\textwidth}\centering
      \colorbox{green!10}{\parbox{0.90\linewidth}{\centering\vspace{0.65em}\textbf{Réponse ancrée}\\articles + JP\vspace{0.65em}}}
    \end{column}
  \end{columns}

  \vspace{0.9em}
  \begin{columns}[T,onlytextwidth]
    \begin{column}{0.47\textwidth}
      \textbf{Sans sources pertinentes}
      \begin{itemize}\setlength{\itemsep}{0.3em}
        \item réponse juridiquement fragile ;
        \item citations inexactes ou inexistantes ;
        \item impossibilité de vérifier le fondement.
      \end{itemize}
    \end{column}
    \begin{column}{0.47\textwidth}
      \textbf{Notre périmètre}
      \begin{itemize}\setlength{\itemsep}{0.3em}
        \item améliorer d'abord le retrieval ;
        \item mesurer ensuite la sélection par le LLM ;
        \item ne pas revendiquer encore une réponse juridique autonome.
      \end{itemize}
    \end{column}
  \end{columns}

  \keybox{\scriptsize En droit, le grounding n'est pas un confort : la source doit être retrouvée,
  identifiable et vérifiable.}
\end{frame}

% 7
\begin{frame}{Que sait-on faire en français sur les textes juridiques et les articles ?}
  \papersection{2. Related Work -- Français et statutory retrieval}
  \scriptsize
  \begin{center}
  \renewcommand{\arraystretch}{1.22}
  \setlength{\tabcolsep}{4pt}
  \begin{tabular}{@{}L{3.0cm} L{4.8cm} L{5.2cm}@{}}
    \toprule
    \textbf{Travail} & \textbf{Apport} & \textbf{Limite pour notre problème} \\
    \midrule
    JuriBERT (2021) & modèle adapté au français juridique & représentation linguistique, pas benchmark de retrieval \\
    BSARD (2022) & questions françaises vers articles belges & articles seulement ; droit belge \\
    Finding the Law (2023) & GNN sur la structure hiérarchique des lois & retrieval statutaire, sans jurisprudence \\
    Open French Law RAG (2024) & étude RAG sur un corpus de droit français & étude de cas, pas benchmark conjoint Articles--JP \\
    \bottomrule
  \end{tabular}
  \end{center}

  \vspace{0.55em}
  \keybox{\scriptsize Le français juridique et le retrieval d'articles disposent déjà de ressources ;
  le lien systématique avec la jurisprudence reste absent de ces travaux.}

  {\tiny Douka et al., NLLP 2021 ; Louis et Spanakis, ACL 2022 ; Louis et al., EACL 2023 ; Harvard LIL, 2024.}
\end{frame}

% 8
\begin{frame}{Que sait-on faire pour retrouver des jurisprudences ?}
  \papersection{2. Related Work -- Legal case retrieval}
  \scriptsize
  \begin{center}
  \renewcommand{\arraystretch}{1.22}
  \setlength{\tabcolsep}{4pt}
  \begin{tabular}{@{}L{2.7cm} L{5.0cm} L{5.3cm}@{}}
    \toprule
    \textbf{Travail} & \textbf{Apport} & \textbf{Limite pour notre problème} \\
    \midrule
    CaseHOLD (2021) & sélection du holding pertinent parmi des candidats & QCM de compréhension ; common law américaine \\
    CaseLink (2024) & graphe inductif pour le retrieval de cas similaires & cas vers cas ; autre système juridique \\
    LePaRD (2024) & retrieval de passages issus de précédents & droit américain ; passages jurisprudentiels uniquement \\
    CLERC (2025) & retrieval de décisions puis génération d'analyse & droit américain ; aucun couple Articles--JP français \\
    \bottomrule
  \end{tabular}
  \end{center}

  \vspace{0.55em}
  \keybox{\scriptsize Le case retrieval est un champ actif : notre différence ne peut pas être
  « retrouver des jurisprudences », mais le contexte français et l'articulation avec les articles.}

  {\tiny Zheng et al., ICAIL 2021 ; Tang et al., SIGIR 2024 ; Mahari et al., ACL 2024 ; Hou et al., NAACL 2025.}
\end{frame}

% 9
\begin{frame}{Que sait-on faire avec des graphes juridiques ?}
  \papersection{2. Related Work -- Graphes et réseaux de citations}
  \scriptsize
  \begin{center}
  \renewcommand{\arraystretch}{1.22}
  \setlength{\tabcolsep}{4pt}
  \begin{tabular}{@{}L{3.0cm} L{4.8cm} L{5.2cm}@{}}
    \toprule
    \textbf{Travail} & \textbf{Apport} & \textbf{Limite pour notre problème} \\
    \midrule
    Finding the Law (2023) & hiérarchie législative intégrée à un GNN & graphe d'articles uniquement \\
    CaseLink (2024) & graphe global pour le retrieval de cas & une seule cible documentaire \\
    LeCNet (2025) & réseau de citations judiciaires et link prediction & droit indien ; prédiction de liens \\
    Belikov et Raoult (2025) & KG de pourvois pénaux de la Cour de cassation & construction de KG, pas benchmark conjoint \\
    \bottomrule
  \end{tabular}
  \end{center}

  \vspace{0.55em}
  \keybox{\scriptsize Les relations juridiques sont déjà exploitées, mais rarement comparées
  comme facteurs contrôlés dans une tâche de retrieval à deux cibles.}

  {\tiny Louis et al., EACL 2023 ; Tang et al., SIGIR 2024 ; Harde et al., JustNLP 2025 ; Belikov et Raoult, 2025.}
\end{frame}

% 10
\begin{frame}{Quel manque précis reste à établir ?}
  \papersection{1--2. Introduction et Related Work -- Gap}
  \footnotesize
  \keybox{\small Les travaux existants étudient séparément la représentation du français juridique,
  le retrieval d'articles, le retrieval de jurisprudences ou la construction de graphes.}

  \vspace{0.7em}
  \begin{columns}[T,onlytextwidth]
    \begin{column}{0.31\textwidth}
      \textbf{\color{accent}Contexte}
      \begin{itemize}\setlength{\itemsep}{0.35em}
        \item droit français ;
        \item sources ouvertes ;
        \item benchmark pénal actuel.
      \end{itemize}
    \end{column}
    \begin{column}{0.31\textwidth}
      \textbf{\color{ok}Deux cibles}
      \begin{itemize}\setlength{\itemsep}{0.35em}
        \item articles applicables ;
        \item jurisprudences pertinentes ;
        \item même question juridique.
      \end{itemize}
    \end{column}
    \begin{column}{0.32\textwidth}
      \textbf{\color{crosscol}Relations évaluées}
      \begin{itemize}\setlength{\itemsep}{0.35em}
        \item citations ;
        \item proximité sémantique ;
        \item architectures hybrides.
      \end{itemize}
    \end{column}
  \end{columns}

  \vspace{0.7em}
  \decisionbox{\scriptsize Gap candidat : aucun benchmark ouvert identifié en droit français
  ne réunit ces trois propriétés. L'affirmation de nouveauté reste à consolider par la bibliographie primaire.}
\end{frame}

% 11
\begin{frame}{Quelle est notre idée principale ?}
  \papersection{1. Introduction -- Key idea et teaser}
  \footnotesize
  \begin{columns}[c,onlytextwidth]
    \begin{column}{0.24\textwidth}
      \centering\textbf{Question $q$}\\[0.4em]
      \small langage naturel
    \end{column}
    \begin{column}{0.05\textwidth}\centering$\Longrightarrow$\end{column}
    \begin{column}{0.36\textwidth}
      \centering
      \colorbox{accentsoft}{\parbox{0.92\linewidth}{\centering\vspace{0.8em}
      \textbf{Graphe juridique typé}\\[0.35em]
      citations + proximités sémantiques\\
      Articles + jurisprudences
      \vspace{0.8em}}}
    \end{column}
    \begin{column}{0.05\textwidth}\centering$\Longrightarrow$\end{column}
    \begin{column}{0.24\textwidth}
      \centering\textbf{Deux classements}\\[0.4em]
      \small sources normatives\\et jurisprudentielles
    \end{column}
  \end{columns}

  \vspace{0.9em}
  \begin{columns}[T,onlytextwidth]
    \begin{column}{0.47\textwidth}
      \textbf{Stage 1 -- Retrieval}\\
      Comparer cosine, PPR et LightGCN sur plusieurs graphes contrôlés.
    \end{column}
    \begin{column}{0.47\textwidth}
      \textbf{Stage 2 -- LLM reranking}\\
      Séparer la couverture du vivier et la capacité du LLM à sélectionner les golds.
    \end{column}
  \end{columns}
\end{frame}

% 12
\begin{frame}{Quelle question scientifique testons-nous ?}
  \papersection{1. Introduction -- Question et portée}
  \footnotesize
  \begin{columns}[T,onlytextwidth]
    \begin{column}{0.66\textwidth}
      \textbf{\color{accent}Question centrale}

      \vspace{0.45em}
      \large
      Dans quelle mesure la structure relationnelle du droit français,
      en particulier les citations entre jurisprudences et articles,
      améliore-t-elle la recherche conjointe des sources normatives et
      jurisprudentielles pertinentes, au-delà de la similarité sémantique seule ?
    \end{column}
    \begin{column}{0.30\textwidth}
      \textbf{Portée empirique}
      \begin{itemize}\setlength{\itemsep}{0.3em}
        \item droit français ;
        \item benchmark pénal ;
        \item Articles et JP séparés ;
        \item aucun score fusionné ;
        \item future lockbox finale.
      \end{itemize}
    \end{column}
  \end{columns}

  \vspace{0.7em}
  \keybox{\scriptsize Le framework vise le droit français ; les conclusions empiriques
  restent limitées au droit pénal français.}
\end{frame}

% 13
\begin{frame}{Que revendiquons-nous exactement ?}
  \papersection{1. Introduction -- Contributions}
  \scriptsize
  \begin{columns}[T,onlytextwidth]
    \begin{column}{0.31\textwidth}
      \textbf{\color{accent}1. Benchmark pénal français}
      \begin{itemize}\setlength{\itemsep}{0.32em}
        \item Dataset d'entraînement et évaluation interne.
        \item Articles et JP retrouvés séparément depuis une même question.
        \item Protocole reproductible.
      \end{itemize}
    \end{column}
    \begin{column}{0.31\textwidth}
      \textbf{\color{ok}2. Architectures de graphes}
      \begin{itemize}\setlength{\itemsep}{0.32em}
        \item Citations et proximités sémantiques.
        \item Graphes citationnels, sémantiques et hybrides.
        \item Relations uniformes, typées ou pondérées.
      \end{itemize}
    \end{column}
    \begin{column}{0.32\textwidth}
      \textbf{\color{crosscol}3. Pipeline retrieval--LLM}
      \begin{itemize}\setlength{\itemsep}{0.32em}
        \item Comparaison de plusieurs retrievers.
        \item Sélection LLM contrôlée sous distracteurs.
        \item Couverture du vivier séparée de l'échec du reranker.
      \end{itemize}
    \end{column}
  \end{columns}

  \vspace{0.8em}
  \keybox{\scriptsize Le papier revendique les ressources, le protocole et l'étude comparative ;
  les conclusions de performance seront formulées après validation des expériences correspondantes.}
\end{frame}

% 8
\begin{frame}{Quelle tâche évaluons-nous ?}
  \papersection{3. Task and Benchmark -- Définition de la tâche}
  \footnotesize
  \begin{columns}[c,onlytextwidth]
    \begin{column}{0.28\textwidth}
      \begin{center}
        \colorbox{accentsoft}{\parbox{0.86\linewidth}{\centering
        \vspace{1.1em}
        \textbf{Question juridique $q$}\\[0.5em]
        énoncée en langage naturel
        \vspace{1.1em}}}
      \end{center}
    \end{column}
    \begin{column}{0.08\textwidth}
      \centering
      $\Longrightarrow$
    \end{column}
    \begin{column}{0.60\textwidth}
      \begin{center}
      \renewcommand{\arraystretch}{1.4}
      \setlength{\tabcolsep}{5pt}
      \begin{tabular}{@{}L{2.9cm} L{4.4cm}@{}}
        \toprule
        \textbf{Sortie} & \textbf{Évaluation} \\
        \midrule
        $R_A^K(q)$ : articles & Recall@10, NDCG@10, MRR@10 \\
        \midrule
        $R_J^K(q)$ : décisions & Hit@10, NDCG@10, MRR@10 \\
        \bottomrule
      \end{tabular}
      \end{center}
    \end{column}
  \end{columns}

  \vspace{1.0em}
  \keybox{\scriptsize Espaces candidats, gold sets, métriques et résultats restent séparés
  par modalité. Aucun score global ne fusionne Articles et JP.}
\end{frame}

% 9
\begin{frame}{Comment le benchmark français est-il construit ?}
  \papersection{3. Task and Benchmark -- Données et protocole de construction}
  \scriptsize
  \begin{center}
  \renewcommand{\arraystretch}{1.32}
  \setlength{\tabcolsep}{4pt}
  \begin{tabular}{@{}c L{2.8cm} L{3.0cm} L{3.0cm} L{3.0cm}@{}}
    \toprule
    & \textbf{Questions + golds} & \textbf{Corpus d'articles} & \textbf{Corpus de décisions} & \textbf{Filtre strict} \\
    \midrule
    \textbf{1} & origine, génération, augmentation & version, période, licence & source ouverte, normalisation & références réellement récupérables \\
    \midrule
    \textbf{2} & extraction des références & identifiants et textes & identifiants, dates et textes & inventaire des exclusions \\
    \bottomrule
  \end{tabular}
  \end{center}

  \vspace{0.7em}
  \textbf{À documenter dans la section :}
  \begin{itemize}\setlength{\itemsep}{0.25em}
    \item statistiques et distribution du nombre de golds ;
    \item provenance, licences et versions des trois composants ;
    \item comptes avant/après filtrage et motifs d'exclusion ;
    \item périmètre pénal, questions générées ou augmentées et gold potentiellement incomplet.
  \end{itemize}

  \decisionbox{\scriptsize À décider : quelles statistiques entrent dans le tableau principal,
  et lesquelles passent au supplément ?}
\end{frame}

% 10
\begin{frame}{Quels graphes comparons-nous ?}
  \papersection{4. Method -- Inventaire du design space $G_0$ à $G_7$}
  \scriptsize
  \begin{center}
  \renewcommand{\arraystretch}{1.12}
  \setlength{\tabcolsep}{3pt}
  \begin{tabular}{@{}L{1.0cm} L{2.7cm} L{4.3cm} L{3.0cm} L{2.0cm}@{}}
    \toprule
    \textbf{ID} & \textbf{Signal} & \textbf{Construction} & \textbf{Rôle} & \textbf{Statut} \\
    \midrule
    $G_0$ & citation brute & graphe JP $\times$ Articles & référence historique & exploratoire \\
    $G_1$ & citation nettoyée & retrait des isolés et non courants & base citationnelle & confirmatoire \\
    $G_2$ & citation sans hubs & $G_1$ moins les plus gros hubs & contrôle hubs & historique \\
    $G_3$ & citation filtrée & $G_2$ moins articles procéduraux & contrôle nettoyage & historique \\
    $G_4$ & sémantique & kNN Article--Article, JP--JP, Article--JP & baseline sémantique & exploratoire \\
    $G_5$ & union naïve & maximum citation + sémantique & contrôle naïf & exploratoire \\
    $G_6/G_{6U}$ & hybride typé & blocs normalisés / contrôle uniforme & structure et contrôle & campagne \\
    $G_7$ & hybride pondéré & poids distincts par famille & régime pondéré & campagne \\
    \bottomrule
  \end{tabular}
  \end{center}

  \vspace{0.45em}
  \keybox{\scriptsize Un second tableau exhaustif liera chaque instance concrète
  à ses paramètres, son \texttt{experiment\_id} et son statut scientifique.}
\end{frame}

% 11
\begin{frame}{Comment le droit est-il représenté ?}
  \papersection{4. Method -- Définition formelle du graphe typé}
  \scriptsize
  \begin{columns}[T,onlytextwidth]
    \begin{column}{0.58\textwidth}
      \centering
      \includegraphics[width=\linewidth,height=0.66\textheight,keepaspectratio]{../specs/graph-visuals-penal-only-2026-06-19/fig_penal_bipartite_sample.png}

      {\tiny Exemple structurel : sous-graphe biparti Articles--JP.}
    \end{column}
    \begin{column}{0.39\textwidth}
      \textbf{Nœuds}
      \begin{itemize}\setlength{\itemsep}{0.25em}
        \item Articles $A$ ;
        \item décisions $J$ ;
        \item attributs textuels et identifiants normalisés.
      \end{itemize}

      \vspace{0.5em}
      \textbf{Relations}
      \begin{itemize}\setlength{\itemsep}{0.25em}
        \item citation Article--JP ;
        \item proximité Article--Article ;
        \item proximité JP--JP ;
        \item proximité Article--JP sémantique.
      \end{itemize}

      \vspace{0.4em}
      \textbf{Paramètres}
      \begin{itemize}\setlength{\itemsep}{0.25em}
        \item poids par famille ;
        \item normalisation par bloc ;
        \item voisinage $k$NN.
      \end{itemize}
    \end{column}
  \end{columns}

  \vspace{0.05em}
  \decisionbox{\tiny À décider : notation finale des matrices par bloc et de leur combinaison.}
\end{frame}

% 12
\begin{frame}{Comment le premier étage classe-t-il les sources ?}
  \papersection{4. Method -- Stage 1 : retrieval}
  \footnotesize
  \begin{columns}[T,onlytextwidth]
    \begin{column}{0.31\textwidth}
      \textbf{\color{accent}Cosine}
      \begin{itemize}\setlength{\itemsep}{0.3em}
        \item Similarité question--document.
        \item Aucun usage du graphe.
        \item Baseline commune à tous les espaces candidats.
      \end{itemize}
    \end{column}
    \begin{column}{0.31\textwidth}
      \textbf{\color{ok}Personalized PageRank}
      \begin{itemize}\setlength{\itemsep}{0.3em}
        \item Seeds issus du signal sémantique.
        \item Propagation non apprise.
        \item Hyperparamètres choisis sur validation.
      \end{itemize}
    \end{column}
    \begin{column}{0.32\textwidth}
      \textbf{\color{crosscol}LightGCN}
      \begin{itemize}\setlength{\itemsep}{0.3em}
        \item Propagation apprise.
        \item Epoch et paramètres sélectionnés sur les folds.
        \item Replay final sans early stopping sur l'eval.
      \end{itemize}
    \end{column}
  \end{columns}

  \vspace{0.9em}
  \keybox{\scriptsize Le negative sampling est une composante ou une ablation de LightGCN,
  pas une quatrième méthode de retrieval.}

  \vspace{0.35em}
  \decisionbox{\scriptsize À décider : statut du contrôle LightGCN non entraîné
  dans le tableau principal.}
\end{frame}

% 13
\begin{frame}{Le LLM sait-il sélectionner les golds sous distracteurs ?}
  \papersection{4--5. Method and Experiments -- Étude LLM contrôlée}
  \scriptsize
  \begin{columns}[T,onlytextwidth]
    \begin{column}{0.58\textwidth}
      \centering
      \includegraphics[width=\linewidth,height=0.64\textheight,keepaspectratio]{../../06-Analyses/comparatifs/llm-rag-rag-k-sweep-2026-07-01/score_vs_rag_k_articles.png}

      {\tiny \color{red!75!black}\textbf{Exploratoire :} forme de courbe servant uniquement à motiver le nouveau contrôle.}
    \end{column}
    \begin{column}{0.39\textwidth}
      \textbf{\color{accent}Contrôle}
      \begin{itemize}\setlength{\itemsep}{0.25em}
        \item Tous les golds restent présents.
        \item Les viviers sont imbriqués.
      \end{itemize}

      \textbf{\color{ok}Difficulté}
      \begin{itemize}\setlength{\itemsep}{0.25em}
        \item Distracteurs issus d'un retriever.
        \item Taille de vivier entre $K=10$ et $K=100$.
      \end{itemize}

      \textbf{\color{warning}Mesure}
      \begin{itemize}\setlength{\itemsep}{0.25em}
        \item Rappel des golds par le LLM.
        \item Rangs et scores masqués.
        \item Ordre randomisé, seeds contrôlés.
      \end{itemize}
    \end{column}
  \end{columns}

  \decisionbox{\tiny À décider : grille exacte de $K$, représentation textuelle,
  budget de contexte et répétitions.}
\end{frame}

% 14
\begin{frame}{Le signal graphe reste-t-il utile après reranking ?}
  \papersection{4--5. Method and Experiments -- Pipeline end-to-end}
  \footnotesize
  \begin{center}
  \renewcommand{\arraystretch}{1.45}
  \setlength{\tabcolsep}{6pt}
  \begin{tabular}{@{}c c c c c c c@{}}
    \colorbox{accentsoft}{\parbox{2.15cm}{\centering\textbf{Question}\\$q$}}
    & $\Longrightarrow$ &
    \colorbox{accentsoft}{\parbox{2.15cm}{\centering\textbf{Retriever}\\cosine ou graphe}}
    & $\Longrightarrow$ &
    \colorbox{accentsoft}{\parbox{2.15cm}{\centering\textbf{Vivier top-$K$}\\sans injection}}
    & $\Longrightarrow$ &
    \colorbox{green!12}{\parbox{1.7cm}{\centering\textbf{Top-10}\\sortie LLM}}
  \end{tabular}
  \end{center}

  \vspace{0.9em}
  \begin{center}
  \renewcommand{\arraystretch}{1.3}
  \setlength{\tabcolsep}{5pt}
  \begin{tabular}{@{}L{4.0cm} L{8.5cm}@{}}
    \toprule
    \textbf{Comparaison} & \textbf{Question répondue} \\
    \midrule
    Retriever seul & Quel top-10 produit le premier étage ? \\
    LLM + baseline cosine & Que gagne le reranking sans graphe ? \\
    LLM + meilleur retriever graphe & Le signal relationnel reste-t-il utile après reranking ? \\
    \bottomrule
  \end{tabular}
  \end{center}

  \vspace{0.55em}
  \keybox{\scriptsize Deux causes d'échec séparées : gold absent du vivier,
  ou gold présent mais non sélectionné par le LLM.}
\end{frame}

% 15
\begin{frame}{Comment rendons-nous les hypothèses falsifiables ?}
  \papersection{5. Experiments -- Protocole confirmatoire}
  \footnotesize
  \begin{center}
  \renewcommand{\arraystretch}{1.4}
  \setlength{\tabcolsep}{5pt}
  \begin{tabular}{@{}L{3.0cm} L{3.2cm} L{3.2cm} L{3.2cm}@{}}
    \toprule
    \textbf{1. Folds groupés} & \textbf{2. Champion gelé} & \textbf{3. Eval interne} & \textbf{4. Lockbox} \\
    \midrule
    train + validation & paramètres et epochs figés & déjà consultée & future et inédite \\
    même provenance dans un fold & aucun choix sur l'eval & confirmation interne & affirmation finale \\
    \bottomrule
  \end{tabular}
  \end{center}

  \vspace{0.8em}
  \textbf{Contrat commun à tous les graphes}
  \begin{itemize}\setlength{\itemsep}{0.3em}
    \item mêmes questions, folds et espaces candidats ;
    \item mêmes seeds, budgets de tuning, métriques et règles de sélection ;
    \item une ablation ne change qu'un facteur causal à la fois ;
    \item tout run non conforme reste explicitement exploratoire.
  \end{itemize}

  \warnbox{\scriptsize \texttt{eval\_rich\_retrievable\_strict} est une confirmation interne,
  jamais une lockbox finale.}
\end{frame}

% 16
\begin{frame}{Quelles comparaisons peuvent soutenir nos affirmations ?}
  \papersection{5. Experiments -- Métriques, baselines et ablations}
  \scriptsize
  \begin{center}
  \renewcommand{\arraystretch}{1.25}
  \setlength{\tabcolsep}{4pt}
  \begin{tabular}{@{}L{2.4cm} L{3.5cm} L{3.5cm} L{3.9cm}@{}}
    \toprule
    & \textbf{Articles} & \textbf{JP} & \textbf{LLM reranking} \\
    \midrule
    Métrique primaire & Recall@10 & Hit@10 & gold recall conditionnel à définir \\
    Métriques secondaires & NDCG@10, MRR@10 & NDCG@10, MRR@10 & top-10 final, perte à l'oracle \\
    Diagnostic & couverture selon $K$ & couverture selon $K$ & robustesse selon taille du vivier \\
    \bottomrule
  \end{tabular}
  \end{center}

  \vspace{0.65em}
  \textbf{Ablations attendues}
  \begin{itemize}\setlength{\itemsep}{0.25em}
    \item citation seule vs sémantique seule ;
    \item union naïve vs blocs typés ;
    \item poids de familles et contrôles uniformes ;
    \item negative sampling au sein de LightGCN ;
    \item retriever seul vs LLM + cosine vs LLM + graphe.
  \end{itemize}

  \decisionbox{\scriptsize À décider : métrique primaire de l'étude LLM contrôlée
  et correction pour comparaisons multiples.}
\end{frame}

% 17
\begin{frame}{Quels tableaux répondront à quelles affirmations ?}
  \papersection{6. Results -- Programme des tableaux}
  \scriptsize
  \begin{center}
  \renewcommand{\arraystretch}{1.2}
  \setlength{\tabcolsep}{4pt}
  \begin{tabular}{@{}L{0.8cm} L{2.8cm} L{5.1cm} L{3.8cm}@{}}
    \toprule
    \textbf{ID} & \textbf{Tableau} & \textbf{Comparaison} & \textbf{Question répondue} \\
    \midrule
    R1 & Retriever principal & cosine / PPR / LightGCN & performance Articles et JP \\
    R2 & Graph design & $G_1$ / contrôles / hybrides & effet des signaux \\
    R3 & Ablations causales & un facteur à la fois & interprétation du gain \\
    R4 & LLM contrôlé & viviers $K=10$ à $100$ & robustesse au contexte \\
    R5 & End-to-end & retriever $\rightarrow$ LLM & top-10 final \\
    \bottomrule
  \end{tabular}
  \end{center}

  \vspace{0.65em}
  \warnbox{\scriptsize Aucun tableau de résultats n'est rempli avant réception
  des exports validés et de leurs \texttt{experiment\_id}.}

  \vspace{0.35em}
  {\scriptsize Les résultats négatifs seront conservés lorsqu'ils expliquent
  une limite ou réfutent une hypothèse du design space.}
\end{frame}

% 21
\begin{frame}{Quelles figures doivent raconter l'argument ?}
  \papersection{6. Results -- Programme visuel}
  \scriptsize
  \begin{columns}[T,onlytextwidth]
    \begin{column}{0.48\textwidth}
      \centering
      \includegraphics[width=\linewidth,height=0.49\textheight,keepaspectratio]{../specs/graph-g0-g1-g2-visuals-2026-06-19/fig_g0_g1_g2_compare.png}

      {\tiny Maquette structurelle : comparaison de familles de graphes.}
    \end{column}
    \begin{column}{0.48\textwidth}
      \centering
      \includegraphics[width=\linewidth,height=0.49\textheight,keepaspectratio]{../../06-Analyses/comparatifs/llm-reranking-depth-2026-07-16/figures/candidate_coverage_articles.png}

      {\tiny \color{red!75!black}\textbf{Exploratoire :} forme de courbe de couverture à reprendre après validation.}
    \end{column}
  \end{columns}

  \vspace{0.3em}
  \textbf{Visuels candidats pour le papier}
  \begin{itemize}\setlength{\itemsep}{0.22em}
    \item pipeline benchmark et définition des deux tâches ;
    \item tableau synthétique $G_0$--$G_7$ et inventaire des instances ;
    \item courbes de couverture selon $K$ ;
    \item courbe de capacité contrôlée du LLM ;
    \item ablations des relations et poids.
  \end{itemize}

  \decisionbox{\tiny À décider : quelles figures entrent dans les 8 pages et lesquelles passent au supplément ?}
\end{frame}

% 22
\begin{frame}{Que pourra-t-on honnêtement conclure ?}
  \papersection{7. Discussion and Limitations}
  \scriptsize
  \begin{center}
  \renewcommand{\arraystretch}{1.25}
  \setlength{\tabcolsep}{5pt}
  \begin{tabular}{@{}L{2.2cm} L{5.0cm} L{2.2cm} L{4.2cm}@{}}
    \toprule
    \textbf{Limite} & \textbf{Conséquence} & \textbf{Limite} & \textbf{Conséquence} \\
    \midrule
    \textcolor{red!75!black}{Causalité} & $G_6$ et $G_7$ changent plusieurs facteurs simultanément & Golds & annotations incomplètes ou clairsemées \\
    Validation & évaluation interne déjà consultée & Portée & droit pénal français uniquement \\
    LLM & sensibilité au prompt, modèle, ordre et contexte & Données & licences, versions et biais de collecte \\
    \bottomrule
  \end{tabular}
  \end{center}

  \vspace{0.8em}
  \textbf{Règles d'interprétation}
  \begin{itemize}\setlength{\itemsep}{0.3em}
    \item une comparaison $G_6$--$G_7$ directe ne mesure pas l'effet isolé d'un seul poids ;
    \item une absence du gold dans le vivier ne doit pas être imputée au reranker ;
    \item la confirmation interne et la future validation lockbox doivent rester distinctes.
  \end{itemize}

  \decisionbox{\scriptsize À discuter : quelles limites doivent apparaître avec le résultat principal,
  plutôt qu'être reléguées en fin de papier ?}
\end{frame}

% 23
\begin{frame}{Que diront la conclusion et l'abstract ?}
  \papersection{8. Conclusion, puis Abstract écrit en dernier}
  \footnotesize
  \begin{columns}[T,onlytextwidth]
    \begin{column}{0.47\textwidth}
      \textbf{\color{accent}Conclusion}
      \begin{itemize}\setlength{\itemsep}{0.45em}
        \item répondre directement à la question de recherche ;
        \item distinguer ce qui est confirmé, réfuté ou encore conditionnel ;
        \item expliquer ce que les résultats changent pour le retrieval juridique ;
        \item ne pas terminer par une liste de travaux non réalisés.
      \end{itemize}
    \end{column}
    \begin{column}{0.47\textwidth}
      \textbf{\color{ok}Abstract -- seulement à la fin}
      \begin{itemize}\setlength{\itemsep}{0.45em}
        \item problème important et encore insuffisamment traité ;
        \item idée principale et entrées/sorties explicites ;
        \item benchmark, méthode et protocole ;
        \item résultats appuyés par leurs \texttt{experiment\_id}.
      \end{itemize}
    \end{column}
  \end{columns}

  \vspace{0.7em}
  \warnbox{\scriptsize Aucun résumé final ni conclusion quantitative ne sera rédigé avant
  réception des preuves confirmatoires validées par la tâche A.}
\end{frame}

% 24
\begin{frame}{Quelles décisions restent à prendre ?}
  {\scriptsize\textcolor{accent}{\textbf{Sortie de la discussion -- Décisions avant rédaction}}}\par\vspace{0.55em}
  \footnotesize
  \begin{center}
  \renewcommand{\arraystretch}{1.35}
  \setlength{\tabcolsep}{5pt}
  \begin{tabular}{@{}c L{11.8cm}@{}}
    \toprule
    \textbf{Ordre} & \textbf{Décision} \\
    \midrule
    \rowcolor{accentsoft}
    \textbf{1} & \textbf{Nom et portée exacte de la contribution méthodologique} \\
    \textbf{2} & Définition formelle du graphe typé, des blocs et des poids \\
    \textbf{3} & Grille $K$, format de contexte et répétitions LLM \\
    \textbf{4} & Baselines, ablations et métriques du reranking \\
    \textbf{5} & Programme final de figures et tableaux \\
    \textbf{6} & Preuves tâche A nécessaires avant rédaction des résultats \\
    \bottomrule
  \end{tabular}
  \end{center}

  \vspace{0.8em}
  \keybox{\scriptsize Ordre du papier : Introduction $\rightarrow$ Related Work $\rightarrow$
  Task and Benchmark $\rightarrow$ Method $\rightarrow$ Experiments $\rightarrow$ Results
  $\rightarrow$ Discussion ; Conclusion et Abstract en dernier.}
\end{frame}

\end{document}
